\documentclass[10pt]{article}
\usepackage[T1]{fontenc}
\usepackage[utf8]{inputenc}
\usepackage[french]{babel}
\usepackage{amsmath,amssymb,mathrsfs,mathpazo}
\usepackage{geometry}
\geometry{a4paper,margin=2cm}
\usepackage{graphicx,array,tabularx,tabularray,colortbl,multirow,multicol,booktabs,enumitem,pifont,ulem,fancybox,diagbox,cancel}
\usepackage{pstricks,pstricks-add,pst-plot,pst-tree,pst-node,pst-text,pst-func,pst-math,pst-bspline,pst-3dplot,pst-eucl,pst-solides3d}
\let\pscalculate\undefined
\usepackage{pst-bezier}
\usepackage{tikz,pgfplots,tkz-tab}
\usetikzlibrary{arrows,arrows.meta,patterns,calc,positioning,trees,decorations.markings,shapes.geometric,shapes.misc,shapes.arrows,intersections,decorations.pathmorphing,decorations.pathreplacing,backgrounds,shadows,babel}
\pgfplotsset{compat=1.18}
\usepgfplotslibrary{fillbetween}
\usepackage[np]{numprint}
\usepackage{listings,makecell,pas-tableur,twoopt}
\let\barre\undefined
\usepackage[pstricks]{tablvar}
\definecolor{darkgreen}{RGB}{0,100,0}
\providecommand{\red}{\color{red}}
\providecommand{\blue}{\color{blue}}
\providecommand{\green}{\color{green}}
\providecommand{\black}{\color{black}}
\providecommand{\white}{\color{white}}
\providecommand{\gray}{\color{gray}}
\providecommand{\cyan}{\color{cyan}}
\providecommand{\magenta}{\color{magenta}}
\providecommand{\yellow}{\color{yellow}}
\providecommand{\violet}{\color{violet}}
\providecommand{\orange}{\color{orange}}
\providecommand{\euro}{\texteuro}
\providecommand{\EUR}{\texteuro}
\providecommand{\ieme}{\textsuperscript{e}}
\providecommand{\degres}{\ensuremath{^\circ}}
\providecommand{\degree}{\ensuremath{^\circ}}
\providecommand{\vv}[1]{\overrightarrow{#1}}
\providecommand{\Coord}[1]{#1}
\providecommand{\decosix}{}
\providecommand{\starredbullet}{\ensuremath{\bullet}}
\pagestyle{empty}
\setlength{\parindent}{0pt}
\usepackage{adjustbox}\DeclareUnicodeCharacter{2236}{:}\DeclareUnicodeCharacter{2190}{\ensuremath{\leftarrow}}
\providecommand{\R}{\mathbb{R}}
\providecommand{\N}{\mathbb{N}}
\providecommand{\Z}{\mathbb{Z}}
\providecommand{\Q}{\mathbb{Q}}
\providecommand{\D}{\mathbb{D}}
\providecommand{\C}{\mathbb{C}}
\providecommand{\e}{{\mathrm{e}}}
\providecommand{\vect}[1]{\overrightarrow{#1}}
\providecommand{\vectt}[1]{\overrightarrow{\mathrm{#1}}}
\providecommand{\vvt}[1]{\overrightarrow{\mathrm{#1}}}
\providecommand{\np}[2][]{{#2~#1}}
\providecommand{\vv}[1]{{\overrightarrow{#1}}}
\renewcommand{\d}{{\mathrm{d}}}
\providecommand{\degres}{{^\circ}}
\providecommand{\degre}{{^\circ}}
\providecommand{\blue}{\color{blue}}
\providecommand{\red}{\color{red}}
\providecommand{\enskip}{\quad}
\providecommand{\llbracket}{\lbrack\!\lbrack}
\providecommand{\rrbracket}{\rbrack\!\rbrack}
\providecommand{\euro}{{\text{€}}}
\providecommand{\no}{{\text{n°}}}
\providecommand{\interval}[2]{{[#1;#2]}}
\providecommand{\intervFO}[2]{{[#1;#2[}}
\providecommand{\dotfill}{\ldots}

\providecommand{\ds}{\displaystyle}
\providecommand{\textding}[1]{#1}
\providecommand{\Oij}{\(\mathrm{O};\vec{\imath},\vec{\jmath}\)}
\providecommand{\Oijk}{\(\mathrm{O};\vec{\imath},\vec{\jmath},\vec{k}\)}
\providecommand{\Ouv}{\(\mathrm{O};\vec{u},\vec{v}\)}
\providecommand{\intervFF}[2]{\ensuremath{[#1;#2]}}
\renewcommand{\euro}{\ifmmode\text{\texteuro}\else\texteuro\fi}
\DeclareUnicodeCharacter{00B0}{\ifmmode^{\circ}\else\textdegree\fi}
\DeclareUnicodeCharacter{20AC}{\ifmmode\text{\texteuro}\else\texteuro\fi}
% Énoncés : Éducation nationale. Transcription et corrigés : APMEP, Denis Vergès et contributeurs. Découpage et classement : CoopMaths, Sébastien Lozano, Éric Elter et contributeurs. Adaptation de la présentation : Maths Explorer.
% CC BY-SA : https://coopmaths.fr/www/about/
% https://coopmaths.fr/alea/static/bac/2026/tex/bac_2026_05_sujet2_ameriquenord_2.tex
\begin{document}
\noindent\textit{APMEP / CoopMaths. CC BY-SA. Présentation adaptée par Maths Explorer.}\par\medskip
\textbf{Partie A: étude du sens de variation d'une fonction}

\medskip

On considère la fonction $f$ définie sur $\R$ par 

\[f(x) = \dfrac{2x}{\sqrt{1 + x^2}}\]

\begin{enumerate}
\item Résoudre l'équation $f(x) = x$.
\item
	\begin{enumerate}
		\item On admet que la fonction $f$ est dérivable sur $\R$.
		
Vérifier que, pour tout réel $x$,on a $f'(x) = \dfrac{2}{(1+x^2)\sqrt{1 + x^2}}$
		\item En déduire le sens de variation de la fonction $f$ sur $\R$.
	\end{enumerate}
\end{enumerate}

\medskip

\textbf{Partie B : étude de la convergence d'une suite récurrente}

\medskip

La suite $(u_n)$ est définie par $u_0 = 1$ et, pour tout entier naturel $n,\: u_{n+1} = f(u_n)$.

\medskip

\begin{enumerate}
\item Démontrer par récurrence que, pour tout entier naturel $n$, on a $1 \leqslant u_n \leqslant u_{n+1} < \sqrt 3$.
\item En déduire que la suite $(u_n)$ converge et déterminer sa limite.
\item Le but de cette question est de retrouver par une autre méthode les résultats de la
question \textbf{2.} de la \textbf{partie B}.

Pour tout entier naturel $n$, on pose :

\[v_n = \dfrac{u_n^2}{3 - u_n^2}\]

On admet que la suite $(v_n)$ est bien définie.
	\begin{enumerate}
		\item Démontrer que la suite $(v_n)$ est une suite géométrique de raison 4 dont on précisera le premier terme.
		\item En déduire une expression de $v_n$ en fonction de $n$ puis que $u_n = \sqrt{\dfrac{1,5 \times 4^n}{1 + 0,5 \times 4^n}}$ pour tout entier naturel $n$.
		\item En déduire la limite de la suite $(u_n)$.
	\end{enumerate}
\end{enumerate}

\medskip

\textbf{Partie C : étude de la convergence de la somme de termes}

\medskip

Pour tout entier naturel $n$ non nul, on pose $S_n = u_0^2 + u_1^2 + \ldots + u_{n-1}^2$.

\medskip

\begin{enumerate}
\item Recopier et compléter le script Python ci-dessous afin que celui-ci permette de lister
les $p$ premiers termes de la suite $(S_n)$.

\begin{center}
\begin{ttfamily}
\begin{tabular}{|l|}\hline
\textbf{from} math \textbf{import}*\\
~\\
\textbf{def} termes(p) :\\
\qquad  u = \ldots\\
\qquad  S=0\\
\qquad  L= [ ]\\
\qquad  \textbf{for} i \textbf{in} range(p)\\
\qquad \qquad S=\ldots\\
\qquad \qquad u= \ldots\\
\qquad \qquad L.append(S) \\
\qquad \textbf{return} L\\ \hline
\end{tabular}
\end{ttfamily}
\end{center}

Remarque : on rappelle qu'en langage Python,

\begin{itemize}[label=$\bullet~$]
\item la commande L= [ ] crée une liste vide ;
\item la commande L.append (S) ajoute, à la fin de la liste L, l'élément supplémentaire S.
\end{itemize}

\item On rappelle que, pour tout entier naturel $k$, on a $1 \leqslant u_k \leqslant \sqrt 3$.

Démontrer que, pour tout entier naturel $n$ non nul, on a $n \leqslant S_n \leqslant 3n$.
\item En déduire les limites respectives de $S_n$ et de $\dfrac{S_n}{n^2}$ lorsque $n$ tend vers $+ \infty$.
\end{enumerate}

\medskip

\end{document}
